\documentclass[10pt,oneside]{report}
\usepackage{soc2style}
\begin{document}
\doctitle{Control Implementation Handbook}{Operational playbooks, safeguards and evidence engineering}{SOC2-03}
\section*{Operating convention and boundaries}
An AICPA criterion identifies an outcome against which controls may be evaluated. The selected point-of-focus headings identify relevant characteristics; they are not discrete certification checklist items. Every owner assignment, technological mechanism and audit test in this suite is a proposed example to be tailored to the system description, commitments, risk profile, period, and service organization responsibilities.

\begin{itemize}
\item Always establish system boundaries, the in-scope trust services categories, service commitments and system requirements first.
\item For a security-only engagement, use the common criteria. For additional trust services categories, include the corresponding additional criteria, subject to properly justified applicability.
\item Document complementary user entity controls and subservice organizations; do not treat vendor assurance reports as replacements for management ownership.
\item Privacy points of focus marked [P] or [C] describe typical processor or controller relevance; assess actual responsibilities rather than treating the annotations as absolute.
\end{itemize}

\chapter{Control design method}
The implementation handbook converts outcomes from the source criteria into proposed technical mechanisms. None of the technologies, thresholds or cadences in this handbook are mandated by the uploaded AICPA publication. Management should choose controls based on the service commitments and system requirements relevant to its SOC 2 engagement.

\subsection{Required design record -- organizational convention, not AICPA template}
\begin{itemize}
\item Control ID and criterion-to-control mapping (many-to-many permitted).
\item Control objective, risk scenario, scope, control owner, performer and reviewer.
\item Control type, triggering event, frequency and dependency on technology general controls.
\item Success condition, failure/exception policy and compensating safeguards.
\item Evidence system, record producer, retention/access controls, population export method and review history.
\item Design approval, change history, walkthrough and periodic revalidation.
\end{itemize}

\chapter{Engineering playbooks}
\section{1. Governance and responsibilities}
Relevant criteria for evaluation: CC1.1 CC1.2 CC1.3 CC1.4 CC1.5 CC5.3. This is a mapping proposal, not a determination of engagement scope.

\begin{engineeringbox}
Define a board-level risk appetite, map each policy to an accountable owner, train personnel, capture exceptions and escalate missed executions.
\end{engineeringbox}
\subsection{Implementation and assurance workflow}
\begin{enumerate}
\item Specify the business commitment, risk and boundaries to be protected; identify dependent control operators and third parties.
\item Write a testable control design including control frequency, accountable owner and expected success criteria.
\item Implement configuration and operational procedures with review and approval, including exception paths.
\item Capture complete source-of-truth populations and event-level records at execution time; restrict evidence modification.
\item Perform an operator walkthrough and independent design challenge; record defects and corrective actions.
\item Review effectiveness at a risk-appropriate interval and preserve evidence of results, not only screenshots of settings.
\end{enumerate}

Suggested evidence package: Quarterly governance packet, signed RACI, policy attestation export and tracked remediation.

Audit challenge: Board oversight is not interchangeable with management approval.

\section{2. System boundary, inventory and data classification}
Relevant criteria for evaluation: CC2.1 CC3.1 C1.1 P3.1. This is a mapping proposal, not a determination of engagement scope.

\begin{engineeringbox}
Discover accounts, endpoints, cloud resources, applications, repositories, personal data stores and outbound integrations. Assign custody, flows, classification and retention.
\end{engineeringbox}
\subsection{Implementation and assurance workflow}
\begin{enumerate}
\item Specify the business commitment, risk and boundaries to be protected; identify dependent control operators and third parties.
\item Write a testable control design including control frequency, accountable owner and expected success criteria.
\item Implement configuration and operational procedures with review and approval, including exception paths.
\item Capture complete source-of-truth populations and event-level records at execution time; restrict evidence modification.
\item Perform an operator walkthrough and independent design challenge; record defects and corrective actions.
\item Review effectiveness at a risk-appropriate interval and preserve evidence of results, not only screenshots of settings.
\end{enumerate}

Suggested evidence package: Inventory snapshots, reconciled cloud configuration exports, system diagram and DFD.

Audit challenge: Separate assets that are provider-managed from systems controlled by the organization.

\section{3. Enterprise risk register and risk acceptance}
Relevant criteria for evaluation: CC3.1 CC3.2 CC3.3 CC3.4 CC5.1. This is a mapping proposal, not a determination of engagement scope.

\begin{engineeringbox}
Link commitments and threats to risk scenarios. Quantify inherent and residual risk using an approved methodology; record treatment, owner, deadline and approval.
\end{engineeringbox}
\subsection{Implementation and assurance workflow}
\begin{enumerate}
\item Specify the business commitment, risk and boundaries to be protected; identify dependent control operators and third parties.
\item Write a testable control design including control frequency, accountable owner and expected success criteria.
\item Implement configuration and operational procedures with review and approval, including exception paths.
\item Capture complete source-of-truth populations and event-level records at execution time; restrict evidence modification.
\item Perform an operator walkthrough and independent design challenge; record defects and corrective actions.
\item Review effectiveness at a risk-appropriate interval and preserve evidence of results, not only screenshots of settings.
\end{enumerate}

Suggested evidence package: Risk register, treatment approvals, material-change reviews and acceptances.

Audit challenge: Do not infer that an arbitrary numerical risk score is required by the TSC.

\section{4. Identity and access governance}
Relevant criteria for evaluation: CC6.1 CC6.2 CC6.3 CC6.6. This is a mapping proposal, not a determination of engagement scope.

\begin{engineeringbox}
Connect HR authoritative records to IAM. Require named human identities, central authentication, privileged role approvals, machine identity inventory and revocation flows.
\end{engineeringbox}
\subsection{Implementation and assurance workflow}
\begin{enumerate}
\item Specify the business commitment, risk and boundaries to be protected; identify dependent control operators and third parties.
\item Write a testable control design including control frequency, accountable owner and expected success criteria.
\item Implement configuration and operational procedures with review and approval, including exception paths.
\item Capture complete source-of-truth populations and event-level records at execution time; restrict evidence modification.
\item Perform an operator walkthrough and independent design challenge; record defects and corrective actions.
\item Review effectiveness at a risk-appropriate interval and preserve evidence of results, not only screenshots of settings.
\end{enumerate}

Suggested evidence package: Identity lists, access requests, role maps, access reviews and removal logs.

Audit challenge: Use independent population joins to detect dormant identities and orphaned service accounts.

\section{5. Privileged access to GitHub and cloud management}
Relevant criteria for evaluation: CC6.1 CC6.3 CC8.1. This is a mapping proposal, not a determination of engagement scope.

\begin{engineeringbox}
Restrict administrative repository/org roles; secure build credentials, OIDC trust, environment approvals, protected branches and privileged operations.
\end{engineeringbox}
\subsection{Implementation and assurance workflow}
\begin{enumerate}
\item Specify the business commitment, risk and boundaries to be protected; identify dependent control operators and third parties.
\item Write a testable control design including control frequency, accountable owner and expected success criteria.
\item Implement configuration and operational procedures with review and approval, including exception paths.
\item Capture complete source-of-truth populations and event-level records at execution time; restrict evidence modification.
\item Perform an operator walkthrough and independent design challenge; record defects and corrective actions.
\item Review effectiveness at a risk-appropriate interval and preserve evidence of results, not only screenshots of settings.
\end{enumerate}

Suggested evidence package: GitHub org audit log, permissions export, cloud IAM policy, repository rulesets and approvals.

Audit challenge: Do not mistake a developer CODEOWNERS approval for platform-wide admin access governance.

\section{6. Secure configuration and vulnerability management}
Relevant criteria for evaluation: CC7.1 CC7.2 CC7.3. This is a mapping proposal, not a determination of engagement scope.

\begin{engineeringbox}
Set hardened baselines, continuously detect drift and known vulnerabilities, deduplicate findings and route to a risk owner with documented SLA and exceptions.
\end{engineeringbox}
\subsection{Implementation and assurance workflow}
\begin{enumerate}
\item Specify the business commitment, risk and boundaries to be protected; identify dependent control operators and third parties.
\item Write a testable control design including control frequency, accountable owner and expected success criteria.
\item Implement configuration and operational procedures with review and approval, including exception paths.
\item Capture complete source-of-truth populations and event-level records at execution time; restrict evidence modification.
\item Perform an operator walkthrough and independent design challenge; record defects and corrective actions.
\item Review effectiveness at a risk-appropriate interval and preserve evidence of results, not only screenshots of settings.
\end{enumerate}

Suggested evidence package: Versioned configuration standards, scan snapshots, alert inventory, remediation ticket timelines.

Audit challenge: A missing scan target must be handled as a coverage failure, not a clean result.

\section{7. Security telemetry and incident evaluation}
Relevant criteria for evaluation: CC7.2 CC7.3 CC7.4. This is a mapping proposal, not a determination of engagement scope.

\begin{engineeringbox}
Define log sources, time synchronization, retention, correlation, on-call triage, severity escalation and incident declaration criteria.
\end{engineeringbox}
\subsection{Implementation and assurance workflow}
\begin{enumerate}
\item Specify the business commitment, risk and boundaries to be protected; identify dependent control operators and third parties.
\item Write a testable control design including control frequency, accountable owner and expected success criteria.
\item Implement configuration and operational procedures with review and approval, including exception paths.
\item Capture complete source-of-truth populations and event-level records at execution time; restrict evidence modification.
\item Perform an operator walkthrough and independent design challenge; record defects and corrective actions.
\item Review effectiveness at a risk-appropriate interval and preserve evidence of results, not only screenshots of settings.
\end{enumerate}

Suggested evidence package: Logging inventory, detection run history, event cases, incident logs and post-incident analyses.

Audit challenge: A detection that never routes to an accountable reviewer is not effective escalation.

\section{8. Incident recovery and resilience}
Relevant criteria for evaluation: CC7.4 CC7.5 CC9.1 A1.3. This is a mapping proposal, not a determination of engagement scope.

\begin{engineeringbox}
Specify recovery objectives, dependencies, restoration approval, integrity verification and full-scale or scoped drills.
\end{engineeringbox}
\subsection{Implementation and assurance workflow}
\begin{enumerate}
\item Specify the business commitment, risk and boundaries to be protected; identify dependent control operators and third parties.
\item Write a testable control design including control frequency, accountable owner and expected success criteria.
\item Implement configuration and operational procedures with review and approval, including exception paths.
\item Capture complete source-of-truth populations and event-level records at execution time; restrict evidence modification.
\item Perform an operator walkthrough and independent design challenge; record defects and corrective actions.
\item Review effectiveness at a risk-appropriate interval and preserve evidence of results, not only screenshots of settings.
\end{enumerate}

Suggested evidence package: Incident runbooks, disaster recovery test packets, validated restore artifacts.

Audit challenge: Failure to validate recovered data may leave a system available but corrupted.

\section{9. Software delivery and change control}
Relevant criteria for evaluation: CC8.1 CC6.3 CC7.1. This is a mapping proposal, not a determination of engagement scope.

\begin{engineeringbox}
Define branch controls and security gates, review changes, require checks before production, secure deployment identity and record exceptions and rollback.
\end{engineeringbox}
\subsection{Implementation and assurance workflow}
\begin{enumerate}
\item Specify the business commitment, risk and boundaries to be protected; identify dependent control operators and third parties.
\item Write a testable control design including control frequency, accountable owner and expected success criteria.
\item Implement configuration and operational procedures with review and approval, including exception paths.
\item Capture complete source-of-truth populations and event-level records at execution time; restrict evidence modification.
\item Perform an operator walkthrough and independent design challenge; record defects and corrective actions.
\item Review effectiveness at a risk-appropriate interval and preserve evidence of results, not only screenshots of settings.
\end{enumerate}

Suggested evidence package: PR approvals, commit SHA, code scanning run metadata, release logs and emergency change approvals.

Audit challenge: Establish which exact revision of code was tested and deployed.

\section{10. SAST and CodeQL assessment}
Relevant criteria for evaluation: CC7.1 CC7.3 CC8.1. This is a mapping proposal, not a determination of engagement scope.

\begin{engineeringbox}
Select in-scope languages and CodeQL query suites, ensure default/protected branch and pull-request analysis, triage findings by risk and verify resolution on the fixed commit.
\end{engineeringbox}
\subsection{Implementation and assurance workflow}
\begin{enumerate}
\item Specify the business commitment, risk and boundaries to be protected; identify dependent control operators and third parties.
\item Write a testable control design including control frequency, accountable owner and expected success criteria.
\item Implement configuration and operational procedures with review and approval, including exception paths.
\item Capture complete source-of-truth populations and event-level records at execution time; restrict evidence modification.
\item Perform an operator walkthrough and independent design challenge; record defects and corrective actions.
\item Review effectiveness at a risk-appropriate interval and preserve evidence of results, not only screenshots of settings.
\end{enumerate}

Suggested evidence package: Code scanning alert snapshots, SARIF uploads, scan status, PR checks and issue history.

Audit challenge: A rule match alone does not establish exploitability; a green workflow is not proof of full repository coverage.

\section{11. Secret scanning and credential exposure}
Relevant criteria for evaluation: CC6.1 CC6.6 CC7.1 CC7.3. This is a mapping proposal, not a determination of engagement scope.

\begin{engineeringbox}
Discover leaked tokens, classify validity and exposure, rotate or revoke applicable credentials, detect misuse and preserve auditable disposition decisions.
\end{engineeringbox}
\subsection{Implementation and assurance workflow}
\begin{enumerate}
\item Specify the business commitment, risk and boundaries to be protected; identify dependent control operators and third parties.
\item Write a testable control design including control frequency, accountable owner and expected success criteria.
\item Implement configuration and operational procedures with review and approval, including exception paths.
\item Capture complete source-of-truth populations and event-level records at execution time; restrict evidence modification.
\item Perform an operator walkthrough and independent design challenge; record defects and corrective actions.
\item Review effectiveness at a risk-appropriate interval and preserve evidence of results, not only screenshots of settings.
\end{enumerate}

Suggested evidence package: Secret scanning alerts, rotation tickets, provider verification and access logs.

Audit challenge: Deleting a secret from Git history without rotation may not neutralize the credential.

\section{12. Open-source dependency governance}
Relevant criteria for evaluation: CC7.1 CC9.2 CC8.1. This is a mapping proposal, not a determination of engagement scope.

\begin{engineeringbox}
Inventory runtime/direct/transitive dependencies, prioritize reachable exposures, verify supported upgrades and validate build/runtime compatibility.
\end{engineeringbox}
\subsection{Implementation and assurance workflow}
\begin{enumerate}
\item Specify the business commitment, risk and boundaries to be protected; identify dependent control operators and third parties.
\item Write a testable control design including control frequency, accountable owner and expected success criteria.
\item Implement configuration and operational procedures with review and approval, including exception paths.
\item Capture complete source-of-truth populations and event-level records at execution time; restrict evidence modification.
\item Perform an operator walkthrough and independent design challenge; record defects and corrective actions.
\item Review effectiveness at a risk-appropriate interval and preserve evidence of results, not only screenshots of settings.
\end{enumerate}

Suggested evidence package: Dependency graph, lockfiles, SBOM, advisory matches and patch PRs.

Audit challenge: CVE existence does not automatically establish exploitability in the deployed application.

\section{13. Third-party and subservice risk}
Relevant criteria for evaluation: CC3.2 CC9.2 CC2.3. This is a mapping proposal, not a determination of engagement scope.

\begin{engineeringbox}
Classify vendor criticality; obtain agreements, service commitments and relevant independent reports; assess carve-outs, complementary user entity controls and unresolved exceptions.
\end{engineeringbox}
\subsection{Implementation and assurance workflow}
\begin{enumerate}
\item Specify the business commitment, risk and boundaries to be protected; identify dependent control operators and third parties.
\item Write a testable control design including control frequency, accountable owner and expected success criteria.
\item Implement configuration and operational procedures with review and approval, including exception paths.
\item Capture complete source-of-truth populations and event-level records at execution time; restrict evidence modification.
\item Perform an operator walkthrough and independent design challenge; record defects and corrective actions.
\item Review effectiveness at a risk-appropriate interval and preserve evidence of results, not only screenshots of settings.
\end{enumerate}

Suggested evidence package: Vendor ledger, due diligence packets, contract clauses and annual review outcomes.

Audit challenge: Verify that the provider report period and relevant services match your system.

\section{14. Confidentiality and privacy data lifecycle}
Relevant criteria for evaluation: CC6.7 C1.1 C1.2 P4.1 P4.2 P4.3. This is a mapping proposal, not a determination of engagement scope.

\begin{engineeringbox}
Maintain labeled data stores, approved uses, encryption and transfer controls, retention clocks, subject-request routing, purge validation and vendor deletion attestations.
\end{engineeringbox}
\subsection{Implementation and assurance workflow}
\begin{enumerate}
\item Specify the business commitment, risk and boundaries to be protected; identify dependent control operators and third parties.
\item Write a testable control design including control frequency, accountable owner and expected success criteria.
\item Implement configuration and operational procedures with review and approval, including exception paths.
\item Capture complete source-of-truth populations and event-level records at execution time; restrict evidence modification.
\item Perform an operator walkthrough and independent design challenge; record defects and corrective actions.
\item Review effectiveness at a risk-appropriate interval and preserve evidence of results, not only screenshots of settings.
\end{enumerate}

Suggested evidence package: Data inventory, retention schedule, data access logs, deletion evidence and handoff records.

Audit challenge: Confidentiality may cover nonpersonal data; do not reduce it to GDPR processing.

\section{15. Availability, backups and capacity}
Relevant criteria for evaluation: CC7.5 A1.1 A1.2 A1.3. This is a mapping proposal, not a determination of engagement scope.

\begin{engineeringbox}
Specify SLO/SLAs, monitor saturation, test scale-up paths, isolate backups, measure RTO and RPO against commitments, and close drill findings.
\end{engineeringbox}
\subsection{Implementation and assurance workflow}
\begin{enumerate}
\item Specify the business commitment, risk and boundaries to be protected; identify dependent control operators and third parties.
\item Write a testable control design including control frequency, accountable owner and expected success criteria.
\item Implement configuration and operational procedures with review and approval, including exception paths.
\item Capture complete source-of-truth populations and event-level records at execution time; restrict evidence modification.
\item Perform an operator walkthrough and independent design challenge; record defects and corrective actions.
\item Review effectiveness at a risk-appropriate interval and preserve evidence of results, not only screenshots of settings.
\end{enumerate}

Suggested evidence package: Capacity forecasts, protected backup inventories, restore metrics and exercise reports.

Audit challenge: A successful backup job does not prove a successful restore.

\section{16. Processing integrity}
Relevant criteria for evaluation: PI1.1 PI1.2 PI1.3 PI1.4 PI1.5. This is a mapping proposal, not a determination of engagement scope.

\begin{engineeringbox}
Document source populations and formats, validate input, control transformation, reconcile outputs and archive records with traceability.
\end{engineeringbox}
\subsection{Implementation and assurance workflow}
\begin{enumerate}
\item Specify the business commitment, risk and boundaries to be protected; identify dependent control operators and third parties.
\item Write a testable control design including control frequency, accountable owner and expected success criteria.
\item Implement configuration and operational procedures with review and approval, including exception paths.
\item Capture complete source-of-truth populations and event-level records at execution time; restrict evidence modification.
\item Perform an operator walkthrough and independent design challenge; record defects and corrective actions.
\item Review effectiveness at a risk-appropriate interval and preserve evidence of results, not only screenshots of settings.
\end{enumerate}

Suggested evidence package: Data contracts, rejected input logs, balancing reports, output manifests and immutable job identifiers.

Audit challenge: Completeness, validity, accuracy, timeliness and authorization must be evaluated separately.

\section{17. Data subject rights and incident notification}
Relevant criteria for evaluation: P1.1 P2.1 P3.2 P5.1 P5.2 P6.1 P6.2 P6.3 P6.6 P6.7 P7.1 P8.1. This is a mapping proposal, not a determination of engagement scope.

\begin{engineeringbox}
Create lawful and role-appropriate privacy workflows for notice, consent, request verification, access, correction, disclosure accounting, complaint handling and breach decision-making.
\end{engineeringbox}
\subsection{Implementation and assurance workflow}
\begin{enumerate}
\item Specify the business commitment, risk and boundaries to be protected; identify dependent control operators and third parties.
\item Write a testable control design including control frequency, accountable owner and expected success criteria.
\item Implement configuration and operational procedures with review and approval, including exception paths.
\item Capture complete source-of-truth populations and event-level records at execution time; restrict evidence modification.
\item Perform an operator walkthrough and independent design challenge; record defects and corrective actions.
\item Review effectiveness at a risk-appropriate interval and preserve evidence of results, not only screenshots of settings.
\end{enumerate}

Suggested evidence package: Published notices, versioned choices, request tickets, disclosure records, legal decisions and reply logs.

Audit challenge: Processor-versus-controller obligations must follow the actual contractual and operating model.

\section{18. Management control monitoring and deficiencies}
Relevant criteria for evaluation: CC4.1 CC4.2 CC5.3. This is a mapping proposal, not a determination of engagement scope.

\begin{engineeringbox}
Define control telemetry and sample-based monitoring; grade failures, escalate overdue items and confirm remediation through retesting.
\end{engineeringbox}
\subsection{Implementation and assurance workflow}
\begin{enumerate}
\item Specify the business commitment, risk and boundaries to be protected; identify dependent control operators and third parties.
\item Write a testable control design including control frequency, accountable owner and expected success criteria.
\item Implement configuration and operational procedures with review and approval, including exception paths.
\item Capture complete source-of-truth populations and event-level records at execution time; restrict evidence modification.
\item Perform an operator walkthrough and independent design challenge; record defects and corrective actions.
\item Review effectiveness at a risk-appropriate interval and preserve evidence of results, not only screenshots of settings.
\end{enumerate}

Suggested evidence package: Control self-assessment, evidence cadence tracker, findings history and retest packets.

Audit challenge: A remedial policy change without evidence of use cannot show operating effectiveness.

\chapter{Practical configuration patterns}
\subsection{Policy-as-code example: change gate design sketch}
This YAML is a specification sketch of a proposed gate. It is not a ready-to-install GitHub ruleset and does not represent AICPA-prescribed workflow syntax. A deployment must translate it into a supported repository rule, required check and branch protection settings.

\begin{lstlisting}[language={},caption={Illustrative release control intent, not GitHub configuration}]
protected_branches: [main, release/*]
required_before_merge:
  - peer_review_approved
  - repository_security_scan_successful
  - tests_passed
  - deployment_credentials_least_privilege
emergency_change:
  requires_incident_or_change_id: true
  post_change_review: true
retained_evidence:
  - commit_sha
  - reviewer_identity_and_timestamp
  - check_run_ids_and_conclusions
  - exception_or_waiver_decision
\end{lstlisting}
\subsection{Security exception record}
\begin{itemize}
\item Exception ID, affected asset or repository, risk scenario, supporting finding and business owner.
\item Reason control cannot operate as designed, likelihood/impact, compensating controls and affected objectives.
\item Proposed expiration, accountable approval independent of the implementer, monitoring requirements and reevaluation triggers.
\item Remediation plan, target date, control retest and explicit closure evidence.
\end{itemize}

\subsection{Minimal evidence manifest}
\begin{lstlisting}[language={},caption={Evidence manifest example (illustrative JSON)}]
{
  "control_id": "APPSEC-CHANGE-01",
  "criterion_ids": ["CC8.1", "CC7.1"],
  "period_start": "YYYY-MM-DD",
  "period_end": "YYYY-MM-DD",
  "population_query": "All production releases in period",
  "exported_at_utc": "YYYY-MM-DDThh:mm:ssZ",
  "source_system": "scm-and-ci",
  "reviewer": "role-not-personal-id",
  "review_result": "complete | exception",
  "immutable_artifact_reference": "approved-evidence-store/key"
}
\end{lstlisting}
\chapter{Control acceptance checklist}
\begin{enumerate}
\item A business objective or commitment can be traced to each in-scope criterion and risk.
\item Management has explicitly approved the control design and assigned an owner.
\item Documented execution has been observed in the actual system, including exceptions.
\item The evidence covers the period and complete in-scope population, with a documented extraction method.
\item Subservice organization dependencies are identified with a documented method of treatment.
\item The control has been reviewed after architecture, risk, personnel and tooling changes.
\item A practitioner can independently reperform a sampled instance without relying only on verbal attestation.
\end{enumerate}


\chapter*{Source and editorial protocol}
\addcontentsline{toc}{chapter}{Source and editorial protocol}
This independent document is grounded in the user-supplied \emph{AICPA TSP Section 100: 2017 Trust Services Criteria for Security, Availability, Processing Integrity, Confidentiality, and Privacy (With Revised Points of Focus---2022)}, 75-page PDF, especially paragraphs .01--.29, the criterion table on pp.~14--67, and the glossary. Criterion wording is quoted when labeled as such. Selected point-of-focus labels are aids to interpretation and may be incomplete. Engineering recommendations, role assignments, tools, tests, dates, scorecards, and example frequencies are author-developed examples and are not AICPA requirements. The 2022 revisions changed points of focus, not the 2017 criteria (source p.~3). The source states that an assessment of each point of focus is not required (p.~5). This document does not constitute a CPA attestation or legal advice.
\end{document}
